\section{Summary}\label{sec:summary}

The proper renormalization of quasi-LF correlations calculated on the lattice has been a main challenge for the applications of LaMET to studying parton physics. Such correlations contain linear divergences associated with Wilson lines which have to be eliminated with high precision to extract the desired physics information. In this work, we propose a self-renormalization method to eliminate both the linear and the logarithmic divergence, as well as the discretization error, from a quasi-LF matrix element which can be matched to a continuum scheme at short distance. As a paradigmatic example, we show that our method works well for $O_{\gamma_t}$ matrix elements in the pion, nucleon, and off-shell quark state. Our analysis shows that the renormalization factors are universal in the hadron state considered. Moreover, the renormalized correlations in the pion state for clover and overlap fermions are similar to each other. Besides, We find a large non-perturbative effect in the popular RI/MOM and ratio renormalization scheme used previously, which has to and can be avoided in the hybrid renormalization scheme proposed recently. It is certainly interesting to see that our method works also for HYP smeared matrix elements. This paves the way to do a phenomenological renormalization for smeared high-precision lattice data. 

For the convenience of the reader, we collect the parameters related to renormalization for various matrix elements discussed in the previous sections, see Tables~\ref{tab:para_hyp0} and~\ref{tab:para_hyp1}. For HYP smearing cases, there is much freedom for the choice of $\Lambda_{\rm QCD}$. We choose $\Lambda_{\rm QCD}=0.39$ GeV in Table~\ref{tab:para_hyp1}. But we also find that $\Lambda_{\rm QCD}=0.1$ GeV gives us smaller fitted discretization errors.


\begin{table}[htbp]
  \centering
  \begin{tabular}{c|cccc}
  \toprule
 Cases & $k$ & $\Lambda_{\rm QCD}$(GeV) & $d$ & $m_{0}$(GeV) \\
\hline
overlap quark & 1.46 & 0.109(02) & -1.29 & 0.0462(16)\\
\hline
clover quark & 1.46 & 0.135(02) & -1.35 & 0.1513(16)\\
\hline
overlap pion & 1.46 & 0.093(10) & -1.17 & -0.0357(46)\\
\hline
clover pion & 1.46 & 0.086(14) & -0.92 &  -0.0715(50)\\
\hline
clover nucleon & 1.46 & 0.093(06) & -0.97 &  -0.0508(40)\\
\hline
  \end{tabular}
  \caption{Renormalization parameters based on the fitting functions Eq.~(\ref{eq:logM}) and Eq.~(\ref{eq:gztoMS}) for the cases without HYP smearing for MILC and RBC ensembles, where $\delta_{\rm sys}=0.002$. $k$ is dimensionless here.}
  \label{tab:para_hyp0}
\end{table}


\begin{table}[htbp]
  \centering
  \begin{tabular}{c|cccc}
  \toprule
 Cases & $k$ & $\Lambda_{\rm QCD}$(GeV)  \\
\hline
overlap quark & 0.5521(07) & 0.39 \\
\hline
clover quark & 0.6328(05) & 0.39 \\
\hline
overlap pion & 0.5191(14) & 0.39 \\
\hline
clover pion & 0.5178(18) & 0.39 \\
\hline
clover nucleon & 0.5139(54) & 0.39 \\
\hline
Wilson Loop   & 0.4921(02) &  0.39 \\
\hline
VEV           & 0.39(21)  &  0.39 \\
\hline
Wilson Link   & \multicolumn{1}{p{1cm}<{\centering}}{0.5402(04), 0.7099(09)} & 0.39  &     &  \\
\hline
  \end{tabular}
  \caption{Renormalization parameters based on the fitting functions Eq.~(\ref{eq:logM_ss}) for the cases with HYP smearing for MILC and RBC ensembles. $k$ is the slope of $e(z)$ (half of the slope for Wilson Loop) and is dimensionless here.}
  \label{tab:para_hyp1}
\end{table}


For comparison purposes, we also show results based on 2+1 flavor clover quark and Luescher-Weisz (equivalent to Symanzik) gauge ensembles from CLS collaboration~\cite{Bruno:2014jqa}. We choose valence fermion action to be clover action, the same as sea fermion action. The renormalization parameters for CLS ensembles are shown in Table~\ref{tab:CLS}. The parameters $k$ and $\Lambda_{\rm QCD}$ for pion matrix elements for CLS ensembles are similar to those for MILC and RBC ensembles, suggesting consistency between results using mixed actions or not. 

\begin{table}[htbp]
  \centering
  \begin{tabular}{c|cccc}
  \toprule
 Cases & $k$ & $\Lambda_{\rm QCD}$(GeV) & $d$ & $m_{0}$(GeV) \\
\hline
clover quark hyp0 & 1.46 & 0.170(06)  &  -0.68 &  0.1936(26) \\
\hline
clover pion hyp0 & 1.46 & 0.113(09)  &  -2.24  &  -0.0941(49) \\
\hline
clover quark hyp1 & 0.573(05)  & 0.39 & 0.53  &  0.2298(26) \\
\hline
clover pion hyp1 & 0.486(10) & 0.39 & -0.21 &  0.0312(29) \\
\hline
  \end{tabular}
  \caption{Renormalization parameters based on the fitting functions Eq.~(\ref{eq:logM}) and Eq.~(\ref{eq:gztoMS}) for CLS ensembles, where $\delta_{\rm sys}=0.002$. “hyp0” is the unsmearing cases and “hyp1” is the HYP smearing cases. $k$ is dimensionless here.}
  \label{tab:CLS}
\end{table}

%In this work, we tested the self-renormalization method for $O_{\gamma_t}$ correlations in pion, nucleon, and off-shell quark states. However, 
We would like to stress that the application of our method is not limited to the matrix elements discussed in this paper. Instead, it can be applied to any matrix element of an operator in the form of Eq.~(\ref{eq:operator}). %We may need to change $Z_{\overline{\mathrm{MS}}}$ (Eq.~(\ref{eq:ZMS})) for different states. 
It can also be generalized, in principle, to any matrix element with a Wilson line
in LaMET applications~\cite{Ji:2020ect}. Of course, in these cases one also needs to consider the physical interpretations of the matrix elements and the details related to lattice calculations carefully.  

 

~\\
~\\
~\\
~\\
~\\
~\\
~\\
~\\
